% ===== Appendix: corrections to the book, what we added, sources (after Questions; not presented) =====
\appendix
\partaccent{lavender}\gdef\idiottopicemoji{books}\gdef\idiottopicname{Հավելված}

\begin{frame}{Հավելված Ա. Ուղղումներ գրքի նկատմամբ (1/3)}
  \footnotesize
  \setlength{\tabcolsep}{4pt}
  \renewcommand{\arraystretch}{1.22}
  \begin{tabular}{@{}L{2.6cm}L{3.2cm}L{4.5cm}L{3.2cm}@{}}
    \toprule
    \textbf{Թեմա} & \textbf{Գրքում (2016)} & \textbf{Այսօր} & \textbf{Աղբյուր}\\
    \midrule
    IQ-ի բաշխումը & 80-110-ում՝ 80\%-ից ավելի & Նորմալ կորով (շեղումը՝ 15)՝ 80-110-ում մոտ 66\%, 80-120-ում՝ 82\% & Վեքսլերի սանդղակներ\\
    «Միջինը միշտ 100 է» & Տրամաբանորեն անհնար է այլ միջին & 100-ը չափաբերման խմբի միջինն է. հին նորմերով այսօրվա միջինը բարձր է & Flynn, Psychol. Bull., 1987\\
    Սպիրմենի g-ն & 1920-ականներ & 1904 & Spearman, Am. J. Psychol., 1904\\
    Բրոկայի գոտին & Գագաթային բլթի գոտի & Ճակատային բլթում է (ստորին ճակատային գալար) & Նեյրոանատոմիայի դասագրքեր\\
    Շարժողական կեղևը & Գագաթային բլթում & Ճակատային բլթի հետին մասն է & Նեյրոանատոմիայի դասագրքեր\\
    \bottomrule
  \end{tabular}
\end{frame}

\begin{frame}{Հավելված Ա. Ուղղումներ գրքի նկատմամբ (2/3)}
  \footnotesize
  \setlength{\tabcolsep}{4pt}
  \renewcommand{\arraystretch}{1.22}
  \begin{tabular}{@{}L{2.6cm}L{3.2cm}L{4.5cm}L{3.2cm}@{}}
    \toprule
    \textbf{Թեմա} & \textbf{Գրքում (2016)} & \textbf{Այսօր} & \textbf{Աղբյուր}\\
    \midrule
    Դանինգ-Կրյուգեր & Ցածր խելք ունեցողները գերագնահատում են իրենց & Խոսքը կոնկրետ հմտության մասին է. պատկերի մի մասը վիճակագրական էֆեկտ է & Nuhfer et al., 2016; Gignac \& Zajenkowski, 2020\\
    Կոռելյացիա 0.2 & Կապը կա 5 մարդուց 1-ի մոտ & Ամբողջ խմբի մասին է. «բացատրում» է տարբերությունների մոտ 4\%-ը & Վիճակագրության հիմունքներ\\
    «25-ում ուղեղը պատրաստ է» & Մեծերի ուղեղը հիմնականում ձևավորված է & Կապերի վերակառուցումը տևում է մինչև մոտ 32 տարեկան & Mousley et al., Nat. Commun., 2025\\
    Ֆլինի էֆեկտ & IQ-ն աճում է ամենուր & 1990-ականներից որոշ երկրներում կանգնել ու նվազել է. պատճառը միջավայրն է & Bratsberg \& Rogeberg, PNAS, 2018\\
    \bottomrule
  \end{tabular}
\end{frame}

\begin{frame}{Հավելված Ա. Ուղղումներ գրքի նկատմամբ (3/3)}
  \footnotesize
  \setlength{\tabcolsep}{4pt}
  \renewcommand{\arraystretch}{1.22}
  \begin{tabular}{@{}L{2.6cm}L{3.2cm}L{4.5cm}L{3.2cm}@{}}
    \toprule
    \textbf{Թեմա} & \textbf{Գրքում (2016)} & \textbf{Այսօր} & \textbf{Աղբյուր}\\
    \midrule
    tDCS & «tCDS». ակնարկը լավացում է գտել & Ճիշտ հապավումը tDCS է. մեկ սեանսի ազդեցությունն առողջների մոտ հուսալի չէ & Horvath et al., Brain Stimul., 2015\\
    Վկաների վստահությունը & «Penrod and Custer» & Cutler և Penrod & Cutler \& Penrod, Mistaken Identification, 1995\\
    Լոլիկի մեջբերումը & «Miles Kingston» & Miles Kington & Kington, The Independent, 2003\\
    Կպելլեների փորձը & 1980-ականներ & Քոուլի խմբի աշխատանքը Լիբերիայում 1960-ականներին էր & Cole et al., The Cultural Context of Learning and Thinking, 1971\\
    \bottomrule
  \end{tabular}
\end{frame}

\begin{frame}{Հավելված Բ. Ինչ ենք ավելացրել գրքից դուրս}
  \begin{columns}[T]
    \begin{column}{0.48\textwidth}
      \begin{card}[colback=mint!30!cardwhite, fontupper=\scriptsize]
        {\blackfont Թվեր ու հետազոտություններ}\par\smallskip
        \begin{itemize}\setlength{\itemsep}{0.2mm}
          \item Դանինգ-Կրյուգերի թվերը՝ Kruger \& Dunning, 1999
          \item Կրթություն ու IQ՝ Ritchie \& Tucker-Drob, 2018
          \item Ուղեղի մարզում՝ Simons et al., 2016; FTC, 2016
          \item Ադերալ առողջների մոտ՝ Ilieva et al., 2013
          \item Ժառանգելիություն ըստ տարիքի՝ Haworth et al., 2010; Plomin \& Deary, 2015
          \item Խելքի գենետիկա՝ Savage et al., 2018
          \item Պիգմալիոն, մտածելակերպ՝ Jussim \& Harber, 2005; Sisk et al., 2018
        \end{itemize}
      \end{card}
    \end{column}
    \begin{column}{0.48\textwidth}
      \begin{card}[colback=peach!35!cardwhite, fontupper=\scriptsize]
        {\blackfont Մերն է, ոչ թե գրքից}\par\smallskip
        \begin{itemize}\setlength{\itemsep}{0.2mm}
          \item «Ծանո՞թ է» ցանկը, «Փորձիր տանը» խորհուրդները
          \item Կենդանիների անունների ու նախշի առաջադրանքները, «Հեղուկ, թե՞ բյուրեղացած» վիկտորինան
          \item Գնահատականների թվերը հորինված են, հասակի ու IQ-ի կետերը՝ մոդելավորված
          \item Ֆրենոլոգիայի նկարը (գլուխ 4-ում չկա)
          \item Բուշի երկրորդ մեջբերումը հանել ենք. աղբյուրը կասկածելի է
          \item Մեմերը. բովանդակությունը գրքից է
        \end{itemize}
      \end{card}
    \end{column}
  \end{columns}
\end{frame}

\begin{frame}{Հավելված Գ. Աղբյուրներ. գլուխ 4}
  \scriptsize
  Գրքի ծանոթագրություններից, որոնց վրա հիմնված են սլայդները\par\smallskip
  \begin{columns}[T]
    \begin{column}{0.48\textwidth}
      \begin{itemize}\setlength{\itemsep}{0.1mm}
        \item Nisbett R. E. et al. Intelligence: new findings... Am. Psychol., 2012
        \item Süß H.-M. et al. Working-memory capacity explains reasoning ability... Intelligence, 2002
        \item Thurstone L. L. Primary Mental Abilities. 1938
        \item Gardner H. Frames of Mind. 1983
        \item Pant A. The astonishingly funny story of Mr McArthur Wheeler. 2014
        \item DeAngelis T. Why we overestimate our competence. APA Monitor, 2003
        \item Rosen H. J. et al. Cognitive self-appraisal in neurodegenerative disease. NeuroImage, 2010
        \item Larson G. E. et al. A «mental effort» hypothesis... Intelligence, 1995
      \end{itemize}
    \end{column}
    \begin{column}{0.48\textwidth}
      \begin{itemize}\setlength{\itemsep}{0.1mm}
        \item Schlaug G. et al. Increased corpus callosum size in musicians. Neuropsychologia, 1995
        \item Maguire E. A. et al. Hippocampi of taxi drivers. PNAS, 2000
        \item Bennabi D. et al. tDCS for memory enhancement. Front. Syst. Neurosci., 2014
        \item Taki Y. et al. Body height, intelligence and gray matter in children. NeuroImage, 2012
        \item Bouchard T. IQ similarity in twins reared apart. 1997
        \item Jerison H. Evolution of the Brain and Intelligence. 1973
        \item Kaino L. M. Traditional knowledge in curricula designs. 2013 (կպելլեներ)
        \item Rosenthal R., Jacobson L. Pygmalion in the classroom. Urban Rev., 1968
      \end{itemize}
    \end{column}
  \end{columns}
\end{frame}
